\documentclass[10pt]{article}
\usepackage[margin=20mm]{geometry}
\usepackage{amsmath,amssymb,mathtools,fontspec,unicode-math,xurl,hyperref}
\IfFontExistsTF{DejaVu Serif}{\setmainfont{DejaVu Serif}}{\setmainfont{Latin Modern Roman}}
\setmathfont{latinmodern-math.otf}
\setlength{\parindent}{0pt}\setlength{\parskip}{6pt}\sloppy
\title{Scientific review of
Tracy–Widom limit for free sum of random matrices}
\author{Initial scientific review: M25-0487}
\date{29 October 2026}
\begin{document}
\maketitle
Abstract: The acquired third preprint contains an invalid moment identity in the final deterministic step of its Green function comparison argument. The separate negativity assertions deduced from that identity are false for admissible absolutely continuous input measures. A concrete two-interval density, paired with a semicircle law, gives a strict sign reversal. The combined coefficient needed by the proof can nevertheless be shown positive by retaining the cancellation between the two measures. The assessment is Major for this core derivation; neither failure of Tracy–Widom universality after correction nor failure of the local laws is established.\par
\section*{Source, hypotheses, and inspected scope}
The paper is by Hong Chang Ji and Jaewhi Park, associated with Annals of Probability 53 (2025), 239–298, DOI 10.1214/24-AOP1705 (https://doi.org/10.1214/24-AOP1705). The reviewed file is the 88-page arXiv:2110.05147v3, submitted 6 March 2023 and dated 8 March on its title page. It contains the orthogonal case and integrated appendices. The journal’s 60-page article and its separate supplement have not been compared. The supplement DOI could not be accessed through the current retrieval route; the acquired manifest also records an unsuccessful supplement retrieval. No identity between the preprint appendices and that supplement is assumed.\par
The model is \(H=A+UBU^*\), with deterministic diagonal \(A,B\) and Haar unitary or orthogonal \(U\). Definition 2.2 and Assumption 2.3 require compact absolutely continuous limiting measures, upper-edge density exponents in \((-1,1)\), Lévy approximation \(N^{-1+o(1)}\), no fixed positive upper outliers, bounded lower endpoints, and zero traces. The measures need not have connected supports. The main result concerns the largest eigenvalue and is accompanied by a joint-edge remark and a local-law theorem.\par
The detailed review follows the theorem hypotheses, the comparison reduction in Section 4, the deterministic coefficient in Section 5.3, and the relevant subordination properties. The lengthy decoupling and fluctuation-averaging estimates in Sections 6–8 and Appendix B have not been independently proved. Reading the orthogonal modifications and the statement of Theorem B.1 is not a certification of all those estimates.\par
\section*{An incorrect sign assertion in the core comparison proof}
Problem 1 (Equations (5.40)–(5.44), pp. 23–24).  For each of the two input measures, the proof asserts \[L_\alpha:=v_\alpha^{\top}w_\alpha
              -m\,v_\alpha^{\top}u_\alpha\le-c<0,
 \qquad
 L_\beta\le-c<0.\] Here \(m=m_{\mu_t}(E_{+,t})\), the vectors \(u_\alpha,v_\alpha\) are defined in (5.28), and \(w_{\alpha k}=\int x^{k-1}(x-\omega_\alpha)^{-3}\,d\mu_\alpha(x)\). The claimed signs are used to bound away from zero the coefficient multiplying \(\operatorname{Im}\mathcal X\) in (5.38).\par
\subsection*{The exact algebra}
Equation reference: eq:L. Suppress the subscript \(\alpha\), let \(\omega>\sup\operatorname{supp}\mu_\alpha\), and put \[m=\int\frac{d\mu_\alpha(x)}{x-\omega}<0,\quad
 q=\int\frac{d\mu_\alpha(x)}{(x-\omega)^2},\quad
 s_-=\int\frac{d\mu_\alpha(x)}{(x-\omega)^3}.\] The definitions on pp. 21–22 give \[v^{\top}u=\frac{q}{m^2}-1,\qquad
 v^{\top}w=m-\frac{2q}{m}+\frac{s_-}{m^2},\] and therefore \[
 L=2m-\frac{3q}{m}+\frac{s_-}{m^2}
   =\frac12F_{\alpha,0}''(\omega)
      +m\bigl(F_{\alpha,0}'(\omega)-1\bigr)
          \bigl(F_{\alpha,0}'(\omega)-2\bigr).\] Equation (5.41) instead gives \(s_-/m^2-q/m\), omitting \(2m-2q/m\). Its last expression has \(mF_{\alpha,0}'(F_{\alpha,0}'-1)\) in place of the second term in [eq:L] (\#eq:L). The omitted term is positive because \(m<0\) and \(q>m^2\). The resulting sign inference is not valid in general.\par
\subsection*{An admissible strict counterexample}
Equation reference: eq:mu. Set \[
 \mu_\alpha=\frac14\operatorname{Unif}[-31/10,-29/10]
             +\frac34\operatorname{Unif}[9/10,11/10],
 \qquad \omega_\alpha=2 .\] This is a compact absolutely continuous probability measure of mean zero. Its upper edge is \(11/10\); the density equals \(15/4\) immediately below that edge, so the required exponent is \(0\).\par
For \(X\sim\mu_\alpha\), put \(R=(2-X)^{-1}\), and denote \(a=\mathbb ER\), \(q=\mathbb ER^2\), \(s=\mathbb ER^3\). Direct integration gives \[\begin{split}
 a&=\frac54\log\frac{51}{49}+\frac{15}{4}\log\frac{11}{9}
       \simeq0.8025217762501908,\\
 q&=\frac{21100}{27489},\qquad
 s=\frac{1739260000}{2266935363}.
 \end{split}\] Here \(m=-a\) and \(s_-=-s\), so [eq:L] (\#eq:L) becomes \[L_\alpha=-\frac{s-3aq+2a^3}{a^2}.\] The retained exact rational calculation verifies \[a^2<q<2a^2,\qquad
 s-3aq+2a^3<-\frac1{25},\qquad
 L_\alpha>\frac7{100}>0 .\] These are strict bounds, not floating-point sign tests. For example, the logarithms were enclosed using \[2\sum_{k=0}^{n-1}\frac{z^{2k+1}}{2k+1}
 \le\log\frac{1+z}{1-z}
 \le2\sum_{k=0}^{n-1}\frac{z^{2k+1}}{2k+1}
       +\frac{2z^{2n+1}}{(2n+1)(1-z^2)}\] with \(n=10\), \(z=1/50\) and \(1/10\). All resulting comparisons use rational arithmetic.\par
It remains to verify that \(\omega_\alpha=2\) is the relevant subordination point for an admissible pair of measures. Take \[b=q^{-1}=\frac{27489}{21100},
 \qquad \mu_\beta=\mu_{\mathrm{sc}}^{(b)} .\] This centered semicircle law has upper-edge exponent \(1/2\). The Pastur equation for \(\mu_\alpha\boxplus\mu_{\mathrm{sc}}^{(b)}\) has its exterior edge at \[E_+=2+ba,\qquad \omega_\beta=a^{-1}+ba .\] Indeed \(m_\alpha(2)=-a=m_\beta(\omega_\beta)\), and \(a^2<q=1/b\) selects the exterior semicircle branch, with \(\omega_\beta>2\sqrt b\). The real map \(\xi\mapsto\xi-bm_\alpha(\xi)\) has derivative \(1-bq(\xi)\). For \(\xi>11/10\), \(q(\xi)\) decreases strictly from infinity to zero, so its unique critical point is \(2\); the second derivative there is \(2bs>0\). Thus this is the regular right edge, not a spurious algebraic branch. Equivalently the two factors in the edge condition (3.9) are \[\frac{q-a^2}{a^2},\qquad \frac{a^2}{q-a^2},\] whose product is \(1\).\par
Assumption 2.3 can also be met by actual deterministic matrix sequences. Use midpoint quantiles of each bounded measure and subtract the empirical mean. Quantile discretization has Kolmogorov error at most \(1/(2N)\); bounded support gives mean error \(O(N^{-1})\), hence Lévy error \(O(N^{-1})\) after centering. The upper endpoints change by \(O(N^{-1})\), the lower endpoints remain bounded, and the traces are exactly zero. Therefore the counterexample to (5.40) satisfies the operative hypotheses, including those needed to realize the limiting measures.\par
\section*{A local repair by combining the two terms}
Equation reference: eq:K. The false individual signs need not be used. At flow time \(t=0\), write \[r=F_{\alpha,0}'(\omega_\alpha)-1>0.\] The edge relation gives \(F_{\beta,0}'(\omega_\beta)-1=r^{-1}\). Also \(q_\alpha=m^2(1+r)\) and \(q_\beta=m^2(1+r^{-1})\). The coefficient in (5.38), denoted here by \(K\), is therefore \[K=-\frac{r^3L_\beta+L_\alpha}{m^4(1+r)^2}.\] Inserting the correct identity [eq:L] (\#eq:L), the terms involving \(m\) cancel: \[
 K=-\frac{F_{\alpha,0}''(\omega_\alpha)
                 +r^3F_{\beta,0}''(\omega_\beta)}
               {2m^4(1+r)^2}>0 .\] The last sign follows from the paper’s valid representation in Lemma A.2: \[F_{\alpha,0}(z)-z
       =\int\frac{d\widehat\mu_{\alpha,0}(x)}{x-z},\] where the positive measure \(\widehat\mu_{\alpha,0}\) has positive mass and support within the input’s endpoint interval. Since \(\omega_\alpha\) lies strictly to its right, \(F_{\alpha,0}''(\omega_\alpha)<0\); the same holds for \(\beta\). The endpoint interval, rather than the possibly disconnected input support itself, suffices.\par
For each fixed admissible model, \(K(0)>0\). The continuity and strict subordination gaps in Lemma 3.3 imply \(K(t)>K(0)/2\) on some small time interval. Since the proof only uses \(0\le t\le N^{-1/3+\chi}\) with sufficiently small \(\chi<1/3\), this interval contains the full comparison flow for large \(N\). This establishes the required local replacement for the deterministic coefficient step, conditional on the preceding decoupling estimates.\par
The dependency is explicit: (5.40)–(5.44) are used to deduce (5.45), completing Proposition 4.7; that proposition supplies the derivative bound for Green function comparison in Proposition 4.3, which feeds Theorem 2.5. The acquired argument is therefore assessed as Major. The verified replacement [eq:K] (\#eq:K) shows why this does not imply that the desired comparison estimate or universality theorem, correctly normalized, is false. It does not certify the unexamined lengthy stochastic expansions.\par
\section*{Normalization and other bounded checks}
Lemma 2.4 defines the edge density coefficient by \(\rho_0(E_+-x)\sim\gamma_0^{3/2}\sqrt{x}/\pi\). Under that convention the normalized statistic is \(\gamma_0N^{2/3}(\lambda_1-E_+)\). The reciprocal printed in Theorem 2.5 and Remark 2.6 on p. 6 is inconsistent with the convention.\par
This distinction has a direct scale test. Replacing \(A,B\) by \(cA,cB\), \(c>0\), preserves the assumptions and changes \(E_+\) to \(cE_+\) and \(\gamma_0\) to \(\gamma_0/c\). The printed statistic consequently changes by \(c^2\), whereas the corrected statistic is unchanged. A universal nondegenerate Tracy–Widom limit cannot be invariant under, for example, multiplication by \(4\). The proof itself provides clear counterevidence to an intended reciprocal convention: (3.13), Definition 4.1, and Proposition 4.2 use \(L_{+,t}=\gamma_tE_{+,t}\), the resolvent of \(\gamma_tH_t\), and the factor \(\gamma_t\). The reciprocal requires correction; no separate severity is based solely on that typesetting inconsistency.\par
The exact identity (4.32) follows from \(\dot E_+=-m_{\mu_t}(E_+)\) and \(\gamma HG=zG+I\). The comparison exponents are consistent: a bound \(N^{5/6+C\epsilon}\) integrated over an energy interval of length \(N^{-2/3+\epsilon}\) and a time interval of length \(N^{-1/3+\chi}\) gives \(N^{-1/6+(C+1)\epsilon+\chi}\), which vanishes when \(\epsilon,\chi\) are chosen sufficiently small. The local-law statement and orthogonal modifications were traced as inputs, without claiming a full proof of their remainder estimates.\par
The findings concern the acquired preprint. The available institutional record links to the preprint, and a bounded correction search did not identify a separate erratum in the retrieved material. The inaccessible journal supplement remains a source limitation.\par
Assessment | Repair verification | Author adoption | Remaining scope\par
Major: (5.40)–(5.44) | Combined coefficient positive; theorem normalization reconciled | Unknown | Full stochastic estimates; journal and supplement comparison\par
\section*{Scope and references}
The retained calculation uses exact rational logarithm bounds and symbolic identities for the moment expansion and the combined coefficient. It contains no author code and no matrix simulation. Exact source, compilation, dependency, and subsequent page-view records are retained with the accompanying evidence. The assigned initial-review AI session performs the scientific analysis and image inspection; the native helper performs only clerical binding checks. No human review or independent acceptance is asserted.\par
9 H. C. Ji and J. Park, Tracy–Widom limit for free sum of random matrices, arXiv:2110.05147v3 (6 March 2023), https://arxiv.org/abs/2110.05147v3. Journal DOI: 10.1214/24-AOP1705. ISTA Research Explorer, publication record 19039, https://research-explorer.ista.ac.at/record/19039, consulted for publication and source-version information.\par
\end{document}
